\documentclass[11pt]{article}
\usepackage[margin=1in]{geometry}
\usepackage{amsmath,amssymb,amsthm}
\usepackage{graphicx}
\usepackage[hidelinks]{hyperref}

\newtheorem{theorem}{Theorem}
\newtheorem{lemma}{Lemma}

\newcommand{\N}{\mathbb N}
\newcommand{\clop}[1]{N_{#1}}
\newcommand{\restr}{\upharpoonright}

\title{A ZFC Continuum-Sized Almost Disjoint Family on \(\mathbb N^*\)}
\author{Run \texttt{fa\_banach\_001}, lane 13}
\date{2026-06-16}

\begin{document}
\maketitle

\section*{Status}

Candidate partial result, likely valid. The packet gives a ZFC lower bound for
Question 4.3 of Lambie-Hanson--Schrittesser \cite{LHS2026}. It does not settle
the full requested cardinalities \(2^{\aleph_1}\) or \(2^{2^{\aleph_0}}\) in
all models.

\section*{Source Question}

The source paper defines an almost disjoint family on a topological space \(X\)
to be a family of noncompact regular open subsets of \(X\) whose pairwise
intersections are compact. Question 4.3 asks:

\begin{quote}
Is it provable in ZFC that there exists an almost disjoint family on
\(\mathbb N^*\) of cardinality \(2^{\aleph_1}\), or even of cardinality
\(2^{2^{\aleph_0}}\)?
\end{quote}

The source crop from page 10 is:

\begin{center}
\includegraphics[width=\textwidth]{figures/open_problem_crop.png}
\end{center}

\section*{Idea of the Proof}

The construction embeds the ordinary almost disjoint-family trick into
\(\mathbb N^*\), but each ordinary almost disjoint set is expanded into a
noncompact regular open set. A branch \(A_x\subseteq\omega\) is split into
infinitely many infinite blocks. The finite initial unions of these blocks give
an increasing sequence of clopen subsets of \(\mathbb N^*\); their union is
open and noncompact. Regularity is the only point to check. If an ultrafilter
is outside the union, then every basic neighborhood of it can be thinned either
outside the branch \(A_x\), or to one point in infinitely many tail blocks, and
the resulting basic clopen set misses the union.

\begin{theorem}
In ZFC, there is an almost disjoint family on \(\mathbb N^*\), in the sense of
Definition 2.1 of \cite{LHS2026}, of cardinality \(2^{\aleph_0}\).
\end{theorem}

\begin{proof}
Identify \(2^{<\omega}\) with \(\omega\). For each branch \(x\in 2^\omega\),
let
\[
  A_x=\{x\restr n:n<\omega\}\subseteq 2^{<\omega}\cong \omega.
\]
Then \(\{A_x:x\in 2^\omega\}\) is a classical almost disjoint family on
\(\omega\) of size \(2^{\aleph_0}\).

For each \(x\in 2^\omega\), fix a partition
\[
  A_x=\bigsqcup_{n<\omega} A_{x,n}
\]
into infinite sets, and for \(n\geq 1\) set
\[
  b_{x,n}=\bigcup_{m<n} A_{x,m}.
\]
For an infinite \(a\subseteq\omega\), write
\[
  N_a=\{p\in\mathbb N^*:a\in p\}
\]
for the corresponding basic clopen subset of \(\mathbb N^*\). Define
\[
  U_x=\bigcup_{n\geq 1}N_{b_{x,n}}.
\]
Each \(U_x\) is open.

We first show that \(U_x\) is regular open. It is enough to prove that every
point outside \(U_x\) is outside the interior of \(\overline{U_x}\). Fix
\(p\in\mathbb N^*\setminus U_x\), and let \(N_y\) be a basic neighborhood of
\(p\), so \(y\in p\) is infinite. Since \(p\notin U_x\), we have
\(b_{x,n}\notin p\) for every \(n\geq 1\). Therefore \(y\setminus b_{x,n}\)
is infinite for every \(n\); otherwise \(y\subseteq^* b_{x,n}\), which would
force \(b_{x,n}\in p\).

If \(y\setminus A_x\) is infinite, choose an infinite
\(z\subseteq y\setminus A_x\). Then \(z\cap b_{x,n}=\emptyset\) for all \(n\),
so \(N_z\subseteq N_y\) is disjoint from \(U_x\).

Otherwise \(y\setminus A_x\) is finite. Since \(y\setminus b_{x,n}\) is
infinite for every \(n\), the set \(y\) meets infinitely many blocks
\(A_{x,k}\). Choose an infinite \(z\subseteq y\cap A_x\) containing at most
one point from each block \(A_{x,k}\). Then \(z\cap b_{x,n}\) is finite for
every \(n\). Hence \(N_z\cap N_{b_{x,n}}=\emptyset\) inside \(\mathbb N^*\)
for every \(n\), and again \(N_z\subseteq N_y\) is disjoint from \(U_x\).

Thus no point of \(\mathbb N^*\setminus U_x\) lies in the interior of
\(\overline{U_x}\). Since \(U_x\) is open, \(U_x=\operatorname{int}
\overline{U_x}\), so \(U_x\) is regular open.

Next, \(U_x\) is not compact. Indeed, it is covered by the increasing open
cover \(\{N_{b_{x,n}}:n\geq 1\}\). No finite subfamily covers \(U_x\), because
\(N_{b_{x,n+1}}\not\subseteq N_{b_{x,n}}\): the infinite set
\(A_{x,n}=b_{x,n+1}\setminus b_{x,n}\) yields nonprincipal ultrafilters in
\(N_{b_{x,n+1}}\setminus N_{b_{x,n}}\). Therefore \(U_x\) is noncompact.

Finally, if \(x\neq y\), then \(A_x\cap A_y\) is finite. Hence
\(b_{x,n}\cap b_{y,m}\) is finite for all \(n,m\), so
\[
  N_{b_{x,n}}\cap N_{b_{y,m}}=\emptyset
\]
inside \(\mathbb N^*\). It follows that \(U_x\cap U_y=\emptyset\), which is
compact. Therefore \(\{U_x:x\in2^\omega\}\) is an almost disjoint family on
\(\mathbb N^*\) of cardinality \(2^{\aleph_0}\).
\end{proof}

\section*{Scope and Limitations}

This is a partial answer to Question 4.3. It proves the ZFC lower bound
\(2^{\aleph_0}\). It does not prove in ZFC that there is such a family of
cardinality \(2^{\aleph_1}\), nor of cardinality \(2^{2^{\aleph_0}}\). In
models satisfying \(2^{\aleph_0}=2^{\aleph_1}\), the theorem reaches the first
cardinal target in Question 4.3, but the general case where
\(2^{\aleph_0}<2^{\aleph_1}\) remains outside this packet.

\section*{Novelty and Search Notes}

The local run indexes were searched for arXiv id 2606.02040, the title, and
core phrases around generalized almost disjoint families on \(\mathbb N^*\).
External web search was performed for the exact source question and close
phrases including ``Ext\({}^2(c_0(\omega_1),c_0)\)'', ``injective dimension
\(c_0\)'', and ``almost disjoint family on \(N^*\) \(2^{\aleph_1}\)''. The
search surfaced the source paper itself but no separate later answer. The
construction above is elementary and should be reviewed as a direct ZFC
partial result rather than as a claimed literature resolution of Question 4.3.

\section*{Human Review Recommendation}

The key point for review is the regular-open verification. In particular,
check the thinning argument for a basic neighborhood \(N_y\) of an ultrafilter
outside \(U_x\). Once regularity is confirmed, noncompactness and pairwise
compact intersections are immediate.

\begin{thebibliography}{9}

\bibitem{LHS2026}
C. Lambie-Hanson and D. Schrittesser,
\emph{Generalized almost disjoint families and injective Banach spaces},
arXiv:2606.02040, 2026.
\url{https://arxiv.org/abs/2606.02040}

\end{thebibliography}

\end{document}

